% !PRINTED 456
\weberpage{456}

If $\beta, \gamma \ldots$ are roots of other equations in $\Omega$, or also other roots of the same equation, then by simultaneous adjunction of $\alpha, \beta, \gamma \ldots$ we likewise obtain algebraic fields over $\Omega$, which we denote by
\[ \Omega(\alpha, \beta, \gamma \ldots) \]
We shall see presently that this extension of the concept of algebraic fields is only an apparent one, and so for the time being we stop at the adjunction of \gesperrt{one} quantity $\alpha$, that is, at the consideration of $\Omega(\alpha)$.

The equation (1), whose root is $\alpha$, may be reducible. But among the irreducible factors of $F(x)$ there is at least one that vanishes for $x = \alpha$, which we shall denote by $f(x)$; this function $f(x)$ is, if we take the coefficient of the highest power of $x$ equal to 1, completely determined, because by the preceding section $x = \alpha$ cannot be the root of two different irreducible equations.

The equation $f(x) = 0$ thus has the form
\weq{2}{x^n + a_1x^{n-1} + a_2x^{n-2} + \cdots + a_{n-1}x + a_n = 0,}
in which $a_1, a_2 \ldots a_n$ are quantities in $\Omega$. If $n$ is the degree of this equation, then we also call $n$ the \gesperrt{degree of the field} $\Omega(\alpha)$.

All quantities $\Theta$ of the field $\Omega(\alpha)$ can be derived by addition, subtraction, multiplication and division from $\alpha$ and from numbers in $\Omega$; they can therefore be represented as rational functions of $\alpha$ with coefficients in $\Omega$, or in the form
\weq{3}{\Theta = \dfrac{\chi(\alpha)}{\psi(\alpha)},}
in which $\chi(x)$ and $\psi(x)$ are functions in $\Omega$. But since $\psi(\alpha)$ must be different from zero, $\psi(x)$ is not divisible by $f(x)$, and since $f(x)$ is assumed irreducible, $\psi(x)$ is relatively prime to $f(x)$.

Accordingly, by \S.\,6, Theorem II, we can determine the functions $Q(x)$, $\varphi(x)$ in $\Omega$, and indeed $\varphi(x)$ at most of degree $n - 1$, so that
\weq{4}{Q(x) f(x) + \varphi(x) \psi(x) = \chi(x)}
holds, and if we then set $x = \alpha$ herein, so that $f(\alpha) = 0$, it follows that
\weq{5}{\dfrac{\chi(\alpha)}{\psi(\alpha)} = \varphi(\alpha).}
